\documentclass[conference]{IEEEtran}
\usepackage[utf8]{inputenc}
\usepackage[T1]{fontenc}
\usepackage{cite}
\usepackage{amsmath,amssymb}
\usepackage{graphicx}
\usepackage{booktabs}
\usepackage{url}

\title{An Autonomous Generate→Verify→Repair (GVR) Pipeline for Intent-to-Configuration Safety in NetOps: A Confirmatory, Paired Evaluation with gpt-5-mini}

\author{
\IEEEauthorblockN{Autonomous AI Researcher}
\IEEEauthorblockA{CNSM 2026 Agentic AI Researcher Track}
}

\begin{document}

\maketitle

\begin{abstract}
We conducted a fully preregistered, autonomous, within-intent paired trial (study-gvr-2026-001) to evaluate whether a generate→verify→repair (GVR) pipeline reduces emulator-confirmed safety violations in LLM-produced network configurations. We sampled 100 public NetOps intents (vendor-stratified) and compared: (A) LLM-only generation (gpt-5-mini) and (B) the same generation guarded by a deterministic verifier plus a single automated repair when the verifier flagged a problem. Deterministic verification covered reachability, ACL semantics/ordering, forwarding-loop detection, and explicit policy contradictions; Mininet-based emulation confirmed operational manifestation. Ninety-nine complete paired intents were analyzed (planned episodes=200, completed=198). Baseline (LLM-only) produced 97/99 successes (no emulator-confirmed violation); the GVR pipeline produced 99/99 successes (paired absolute improvement = 0.02020202). Exact McNemar two-sided on discordant pairs (n\_10=2, n\_01=0) yields p = 0.5; paired bootstrap-percentile 95\% CI (10k resamples, seed=1729) = [0.000, 0.05050505]. The preregistered large-effect target (>=30 percentage-point absolute reduction) was not met. We describe failure modes, operational implications, and release the full reproducibility artifact bundle per the Disclosure Statement.
\end{abstract}

\section{Introduction}
Motivation and question. Large language models (LLMs) offer a promising path for intent-driven NetOps automation (natural-language intent → device CLI). However, stochastic generation risks producing misconfigurations that can violate reachability, policy, or availability. Generate→Verify→Repair (GVR) architectures interpose deterministic validators (oracles) and repair loops between probabilistic generation and acceptance; prior conceptual and single-case work suggests they can increase operational safety but rigorous, preregistered empirical evidence for NetOps is limited. We therefore ask: Does an end-to-end GVR pipeline (LLM + deterministic network verifier + emulator confirmation) materially reduce emulator-confirmed safety violations compared to LLM-only generation for a public NetOps intent corpus? Contributions. (1) A preregistered, autonomous, within-intent paired trial (study-gvr-2026-001) comparing LLM-only versus a GVR pipeline using gpt-5-mini; (2) deterministic, auditable evidence (prompts, raw outputs, verifier logs, emulator logs, provider-call traces, analysis code) released with content-addressed hashes; (3) a focused analysis that explains why a pre-registered large-effect threshold was not observed and outlines operational implications for GVR adoption.

\section{Related Work}
Positioning. Our work intersects (i) intent-to-configuration research that evaluates LLM feasibility for NetOps translation, (ii) deterministic network-verification systems (e.g., VeriFlow/NetKAT family) that provide oracle-style invariants for reachability and policy, and (iii) generate-verify-repair and convergent-correctness architectures used in regulated-code contexts. Surveys of retrieval- and verifier-augmented LLMs report grounding benefits but heterogeneous evaluation practices; we contribute a preregistered, emulator-confirmed paired comparison focusing on operational failure manifestation rather than token-level correctness alone. Representative verified and RAG work informed our verifier design and failure-mode taxonomy.

\section{Methodology}
Preregistration and experimental plan. The full preregistration (study-gvr-2026-001) froze design choices before any model calls: N=100 public English NetOps intents stratified by vendor-family (Cisco=34, Juniper=33, Arista=33), within-intent paired design, single shared initial generation per intent, maximum one automated repair call per intent (guarded condition), deterministic retry policy, contamination/missingness rules, and primary/secondary estimands. Pipeline implementation. LLM: gpt-5-mini (hosted API). Generation: one shared-initial candidate per intent (one model call) was produced and used for both conditions. Guarded (GVR) path: the shared candidate was fed to a deterministic verifier implementing four predicate classes: (1) reachability (source→destination connectivity and missing routes), (2) ACL semantics \& ordering contradictions (deny/allow conflicts and placement effects), (3) forwarding-loop static detection, and (4) explicit intent-policy contradictions. If the verifier flagged any predicate the pipeline invoked a single repair call that injected typed verifier diagnostics into a repair prompt; at most one repair call per intent was allowed per preregistration. Emulation confirmation. We required an "emulator-confirmed safety violation" to be (a) a verifier flag AND (b) Mininet-based emulation reproducing the operational failure when probed (ICMP/TCP probes and flow observations). Scoring and analysis. Primary estimand: paired absolute difference in proportions of intents producing >=1 verifier-flagged, emulator-confirmed violation (guarded minus baseline). Primary test: exact McNemar two-sided on discordant pairs; paired bootstrap-percentile 95\% CIs (10,000 resamples, seed=1729). Reproducibility. All orchestration, seeds, prompts, raw LLM responses, provider-call traces, verifier JSON reports, emulator logs, and analysis code were archived and content-addressed; critical artifact hashes appear in the Disclosure Statement.

\section{Results}
Execution summary. Planned episodes = 200 (100 intents × 2 conditions). Completed episodes = 198; failed episodes = 2 (one baseline, one guarded) handled per preregistered retry rules; final paired analysis used 99 complete intent pairs. Primary outcomes. Baseline (LLM-only) produced 97/99 successes (no emulator-confirmed violation); GVR (guarded) produced 99/99 successes. Paired contingency: n\_11 = 97 (success both), n\_10 = 2 (baseline fail → guarded success), n\_01 = 0 (baseline success → guarded fail), n\_00 = 0. Point estimate: paired absolute increase in success proportion (guarded - baseline) = 0.02020202 (2.02 percentage points). Exact McNemar two-sided on discordant pairs (n=2) yields p = 0.5. Paired bootstrap-percentile 95\% CI (10k resamples, seed=1729) for the difference = [0.000, 0.05050505]. Provider and artifact accounting. The execution manifest reports model\_calls\_used = 102 (shared initial = 100, repair = 2); full provider-call traces and scoring JSONs are archived for inspection. Missingness and contamination. The run followed the preregistered contamination and missingness procedures; contamination flags and per-intent details are published in the artifact bundle.

\section{Discussion}
Interpretation. The preregistered large-effect target (>=30 percentage-point absolute reduction) was not achieved. Three operational explanations account for the near-null confirmatory outcome: (1) observed baseline correctness was unexpectedly high (97/99 successes), removing headroom for large absolute improvements under the pre-registered sample size and effect-size assumption; (2) verifier scope was intentionally conservative (limited to reachability, ACL/order, forwarding-loop, explicit-policy predicates) and therefore missed other fault classes; (3) the repair budget (one automated repair call per intent) constrained iterative convergence on harder faults; some repairs require multiple generate→verify→repair iterations. Statistical implications. The preregistered power calculation assumed a baseline failure rate ≈40\% to detect a 30pp drop with N=100; with the realized baseline ≈2\% the original power targets do not apply and the trial was not expected to observe large absolute reductions. Operational implications. Deterministic verification remains valuable as an audit and diagnostic layer and can convert some failing LLM outputs into safe configurations (2 observed conversions in this run). To realize measurable large gains in practice, organizations should consider (a) targeting harder distributions (lower baseline correctness), (b) expanding verifier predicates and fidelity, (c) permitting multi-iteration repair budgets, or (d) combining ensemble/multi-candidate generation with verifier-guided selection. Reproducibility and artifacts. We publish the complete artifact bundle (prompts, raw responses, verifier reports, emulator logs, provider-call traces, and deterministic analysis environment) to enable independent reproduction and alternative sensitivity analyses; critical hashes are listed in the Disclosure Statement. Limitations. Key threats include single-model evaluation (gpt-5-mini), constrained repair budget, conservative verifier coverage, and Mininet emulation fidelity; see Limitations section for details.

\section{Limitations}
\begin{itemize}
\item Single-model evaluation and external validity: the experiment used a single hosted LLM (gpt-5-mini) and a 100-intent public sample; results may not generalize to other models, private operational intent distributions, or larger, more diverse production corpora.
\item Verifier coverage and oracle limits: verifier predicates were intentionally conservative (reachability, ACL/order, forwarding-loop, explicit-policy contradictions). Many operational correctness dimensions (timing behaviors, vendor micro-behaviors beyond encoded rules) are outside verifier scope.
\item Repair budget constraint: the preregistered maximum of one repair call per intent limited iterative convergence on harder faults and therefore constrained observable GVR effectiveness.
\item Emulation fidelity: Mininet-based emulation approximates device behavior but does not fully reproduce production hardware, device-specific quirks, or operator workflows; emulator-confirmed failure is a conservative but imperfect proxy for production risk.
\item High baseline correctness and statistical power: the sampled public intents produced a very high baseline no-violation rate (97/99), reducing headroom for measurable improvements relative to preregistered power assumptions (which assumed a much higher baseline failure rate).
\item Single-oracle dependence: deterministic verification combined with a single LLM may miss complementary gains available from multi-model ensembles, richer retrieval/context grounding, or human-in-the-loop approval workflows.
\end{itemize}

\section{Conclusion}
We executed a fully preregistered autonomous paired trial evaluating a GVR pipeline for intent-to-configuration safety. Under the preregistered single-repair budget and conservative verifier predicates, GVR produced a small absolute improvement (2.02 percentage points) on a public intent corpus with high baseline correctness; the preregistered target (>=30pp reduction) was not met. Practical inference: deterministic verification provides precise diagnostics and can fix some faults, but measurable large improvements require either harder tasks, richer verifiers and emulators, iterative repair budgets, or ensemble grounding. All artifacts and analysis code have been released with content-addressed hashes to support independent reproduction and follow-up studies.

\section*{Disclosure Statement}
Mandatory disclosure (autonomy, tools, and process): This study was executed under the CNSM Agentic AI Researcher track preregistration (study-gvr-2026-001). LLM: gpt-5-mini (hosted API). Exact master prompt family, prompt templates, and all concrete prompt strings used for initial generation and repair are archived in the execution bundle (execution/raw\_results.jsonl; hash: 5b7b5cc9ad6529fbb1fc8a7861247d1d53b0779c776d1e081eb1fdce136e3214). Orchestration framework: a deterministic GVR adapter (adapter\_family = hosted\_netops\_gvr\_v1) implemented the generate→verify→repair pipeline and enforced the preregistered limits (shared initial candidate; maximum 1 repair call per intent). Topic constraint: a human Research Director selected the broad topic family (Generative AI / LLMs for NetOps) before the autonomy lock; the autonomous researcher performed all subsequent literature review, study selection, preregistration, experiment execution, analysis, manuscript drafting, peer-review response, and revision after the autonomy lock. Code/experiment-generation framework: Python orchestration executed in Docker with deterministic seeds; orchestration code, Docker image, and deterministic seeds are archived (execution/model\_configuration.json hash: 1acce92558edeb13cd97b75817329d332afe4a36d01d566f1a26c61b132923fa). Preregistration/freeze: the full preregistration (study-gvr-2026-001) preceded any model calls and is unchanged; execution adhered to the preregistered sampling, retry, and stopping rules. Autonomous literature procedure, experiment execution, analysis procedure, manuscript generation, and autonomous peer-review and revision were performed by the autonomous system after the lock; no human manually edited technical manuscript content after the autonomy lock. Reference verification: cited records were programmatically retrieved and metadata-verified; verification records are archived. Technical restarts: two infrastructure failures occurred during runtime and were handled per the preregistered retry policy; these events are logged in execution/execution\_log.jsonl (hash: 0aa7e9e73e551f3625fb2b54af4d0611f302f002606957437ad2800465add271); no post-hoc changes to experimental design were made in response to observed outcomes. Selected artifact pointers and critical hashes: execution/raw\_results.jsonl (5b7b5cc9ad6529fbb1fc8a7861247d1d53b0779c776d1e081eb1fdce136e3214), analysis/paired\_contingency\_table.csv (7cd45c3bd39abae3e068ee546193c2569487071266e3d23095b8b65b2b5e7bcf), analysis/contamination\_summary.csv (389227f0d1e81239c7498abd4944a04c59c92c5fa042b990e332796cd61c3222), analysis/missingness\_summary.csv (764e5d1b87fdf00afbf6c51050b62b39d2d9f58fb033e652dc7e268aa00daa61), execution/provider\_calls/task-000013-guarded-repair.json (dbdfaacf4a3316b490c0ed86d49b2f3d98c317d93d16a816577195070342e45d), execution/provider\_calls/task-000088-guarded-repair.json (9b34fa6be000c542900bd7bc13a9cc020a0a43e66f980d4d673c4d465eced611), and analysis/analysis\_log.jsonl (a6efe50e807252ce2eb7b484b0d5efca52c8cc089147c80356c67ac0905c500c). The canonical confirmatory analysis used the preregistered exact McNemar two-sided conditional test and paired bootstrap-percentile CIs (10,000 resamples, seed=1729) as implemented in the archived analysis container. The autonomous run produced 198 completed episodes (100 shared initial generations and 98 guarded/baseline pairs completed as planned with two infrastructure failures); complete\_pair\_count = 99 after applying preregistered exclusion/missingness rules. All artifacts required to reproduce every reported number are included in the released execution bundle; hashes above permit integrity checking. Human role: humans only set the pre-lock broad topic and constraints; all subsequent research actions qualified by the track were autonomous.

\begin{thebibliography}{99}
\bibitem{ref1} Rafael Cadenas, "Convergent Correctness in Stochastic Code Generation: A Generate-Verify-Repair Architecture for Deterministic Validation of Autonomous AI Coding Agents in Regulated Environments", 2026, 10.2139/ssrn.6754899
\bibitem{ref2} Nguyen Tu, Sukhyun Nam, James Won‐Ki Hong, "Intent‐Based Network Configuration Using Large Language Models", 2024, 10.1002/nem.2313
\bibitem{ref3} Ahmed Khurshid, Xuan Kelvin Zou, Wenxuan Zhou, Matthew Caesar, P. Brighten Godfrey, "VeriFlow: verifying network-wide invariants in real time", 2013
\bibitem{ref4} Raghubir Singh, Sukhpal Singh Gill, "Edge AI: A survey", 2023, 10.1016/j.iotcps.2023.02.004
\bibitem{ref5} Jan Clusmann, Fiona R. Kolbinger, Hannah Sophie Muti, Zunamys I. Carrero, Jan‐Niklas Eckardt, Narmin Ghaffari Laleh, Chiara Maria Lavinia Löffler, Sophie-Caroline Schwarzkopf, Michaela Unger, Gregory Patrick Veldhuizen, Sophia J. Wagner, Jakob Nikolas Kather, "The future landscape of large language models in medicine", 2023, 10.1038/s43856-023-00370-1
\end{thebibliography}

\end{document}
